雅化的利润修复尚未完成现金质量的验证。2025 年归母净利润恢复至 6.32 亿元，但经营活动现金流为 -5.70 亿元；2026Q1 归母为 3.39 亿元而经营现金流为 -4.31 亿元。对周期公司而言，利润增长若主要被应收和库存吸收，低 PE 可以是现金折价而不是低估。

\section{盈利修复与现金背离同时存在}

\begin{exhibitbox}[表 5-1：历史财务与现金转换]
\begin{tabularx}{\exhibitboxwidth}{L{3.4cm}R{2.0cm}R{2.0cm}R{2.0cm}R{2.0cm}}
\toprule
人民币亿元 & 2023A & 2024A & 2025A & 2026Q1 \\
\midrule
营业收入 & 118.95 & 77.16 & 85.43 & 28.30 \\
归母净利润 & 0.40 & 2.57 & 6.32 & 3.39 \\
扣非归母净利润 & -2.14 & 1.64 & 6.94 & 3.57 \\
经营活动现金流 & 8.31 & 9.44 & -5.70 & -4.31 \\
经营现金流 / 归母 & 20.7 倍 & 3.7 倍 & -90.1\% & -127.1\% \\
期末存货 & 22.31 & 16.45 & 16.96 & 18.73 \\
总负债 & 38.12 & 32.60 & 38.99 & 37.65 \\
资本开支代理项 & 6.21 & 5.59 & 3.13 & 0.56 \\
\bottomrule
\end{tabularx}
\end{exhibitbox}
\sourcenote{公司 2023--2025 年年度报告、2026 年第一季度报告；经营现金流/归母为 AStock 复算。}

2023--2024 年的高现金转换包含低盈利基数带来的比例失真，不能据此认为现金风险不存在；2025 年和 2026Q1 的负经营现金流则必须在估值中被视为真实约束。我们的基准情景只假设 2026 年经营现金流约 2 亿元，明显低于 2026 年归母增幅所暗示的会计利润改善，因此目标价没有按高盈利情景给予完全的现金流倍数。

\section{应收与存货：H1 的第一道证伪线}

2026Q1 末应收账款为 19.80 亿元，较年初增加 6.55 亿元；存货为 18.73 亿元，较年初增加 1.77 亿元。它们不能单独证明利润质量恶化，因为锂盐和项目业务具有季节性、库存和结算节奏，但与负经营现金流并列时，已足以要求 H1 提供更完整解释。

\begin{exhibitbox}[表 5-2：正式 H1 的现金质量评分卡]
\begin{tabularx}{\exhibitboxwidth}{L{3.4cm}X X}
\toprule
验证项 & 支持基准/上修的信号 & 触发降级的信号 \\
\midrule
经营现金流与归母 & 现金流转正或有可复算的季节性/结算解释 & 利润上升但现金继续明显为负且无解释 \\
应收账款 & 增长与收入、项目结算和回款计划匹配 & 应收增速持续高于收入且账龄恶化 \\
存货与价格 & 存货规模、跌价准备和价格变化可相互解释 & 存货继续占用且利润依赖未验证的库存收益 \\
分部盈利 & 锂盐与民爆收入、毛利方向支持情景 & 合并利润与分部收入/毛利逻辑不匹配 \\
\bottomrule
\end{tabularx}
\end{exhibitbox}

\section{资产负债表并非当前主风险，但决定估值的现金折扣}

2026Q1 末总资产为 150.75 亿元、负债为 37.65 亿元，杠杆水平本身不是本报告的主要下行来源。风险在于营运资本消耗与扩产/资源安排对现金的竞争：如果利润修复不能转为经营现金流，周期高点的 PE 不能被直接赋予现金复利公司的估值倍数。第八章因此采用周期校准的 EPS 倍数，并用市场锚限制乐观估值的权重。
